% Краткое сообщение ученице:
% Екатерина, в работе осталось несколько точечных ошибок. В третьем задании при вычислении угла не был учтён ещё один известный угол. В шестом итоговые ответы записаны верно, но промежуточное вычисление противоречит условию о медиане. В восьмом расчёты приводят к третьей диаграмме, хотя в ответе указана четвёртая. В десятом неверно оценено первое утверждение. В одиннадцатом и во второй части шестнадцатого решение не завершено. В восемнадцатом неверно выполнено деление при нахождении среднего балла, а в двадцатом найден соседний угол вместо требуемого.
% Репозиторий: students/ekaterina/review_docs/ekaterina_review_2026-09-27.tex
\documentclass[12pt,a4paper]{article}
\usepackage[a4paper,margin=2cm]{geometry}
\usepackage[T2A]{fontenc}
\usepackage[utf8]{inputenc}
\usepackage[russian]{babel}
\usepackage{amsmath,amssymb,mathtools}
\usepackage{array,booktabs,longtable}
\usepackage{xcolor}
\usepackage[unicode,colorlinks=true,linkcolor=blue!55!black,urlcolor=blue!55!black]{hyperref}
\usepackage{tikz}

\setlength{\parindent}{0pt}
\setlength{\parskip}{0.55em}
\setlength{\emergencystretch}{2em}
\raggedbottom

\newcommand{\errorlabel}{\textcolor{red}{\textbf{Ошибка:}}}
\newcommand{\keyidea}{\textcolor{blue!60!black}{\textbf{Ключевая идея:}}}
\newcommand{\resultlabel}{\textcolor{teal!60!black}{\textbf{Итог:}}}
\newcommand{\subhead}[1]{\par\medskip\textbf{#1}\par}

\makeatletter
\def\ps@reviewfooter{%
  \def\@oddhead{}\def\@evenhead{}%
  \def\@oddfoot{\hfil\footnotesize Лёвин Артём Александрович эксклюзивно для Екатерины Скелли\hfil}%
  \let\@evenfoot\@oddfoot
}
\makeatother

\begin{document}

\begin{titlepage}
\thispagestyle{empty}
\centering
\vspace*{0.25\textheight}
{\LARGE\bfseries Ревью домашней работы\par}
\vspace{2em}
{\large Автор: Лёвин Артём Александрович\par}
\vspace{1em}
{\large 27 сентября 2026 года\par}
\vfill
\begin{tikzpicture}[scale=1]
  \draw[line width=0.6pt] (-4,0) .. controls (-2,0.55) and (2,-0.55) .. (4,0);
\end{tikzpicture}
\end{titlepage}

\pagestyle{reviewfooter}

\newpage
\section*{Сводная таблица}
\small
\setlength{\tabcolsep}{3pt}
\begin{longtable}{@{}>{\centering\arraybackslash}p{0.07\linewidth}p{0.16\linewidth}p{0.30\linewidth}p{0.19\linewidth}p{0.20\linewidth}@{}}
\toprule
\textbf{№} & \textbf{Тема} & \textbf{Суть ошибки} & \textbf{Ваш ответ} & \textbf{Правильный ответ}\\
\midrule
\endfirsthead
\toprule
\textbf{№} & \textbf{Тема} & \textbf{Суть ошибки} & \textbf{Ваш ответ} & \textbf{Правильный ответ}\\
\midrule
\endhead
\hyperref[task:3]{3} & Параллельные прямые и углы & При вычислении не учтён угол $22^\circ$; из развёрнутого угла вычтен только угол $72^\circ$. & $108^\circ$ & $86^\circ$\\
\addlinespace
\hyperref[task:6]{6} & Среднее, медиана, диаграмма & Итоговые ответы верны, но промежуточный набор значений даёт медиану $177$, а не $180$. & а) 3; б) 180 ц/га & а) 3; б) 180 ц/га\\
\addlinespace
\hyperref[task:8]{8} & Среднее, медиана, доли & Полученные значения соответствуют диаграмме 3, но в ответе указана диаграмма 4. & 1) 4; 2) 2700 тыс. км$^2$ & 1) 3; 2) 2700 тыс. км$^2$\\
\addlinespace
\hyperref[task:10]{10} & Окружность и геометрические утверждения & Утверждение 1 ошибочно признано верным: для внешнего касания сравнивают расстояние между центрами с суммой радиусов, а не диаметров. & Верны 1, 3, 4 & Верны 3, 4\\
\addlinespace
\hyperref[task:11]{11} & Граф и таблица дорог & Решение и числовой ответ не даны. & «Не знаю» & 12 км\\
\addlinespace
\hyperref[task:16]{16} & Средняя стоимость смеси & Первая часть решена верно, вторая часть не доведена до уравнения и ответа. & 1) 262 руб.; 2) не завершено & 1) 262 руб.; 2) 0,8 кг\\
\addlinespace
\hyperref[task:18]{18} & Среднее арифметическое по диаграмме & Сумма баллов $81$ найдена верно, но $81:22$ вычислено как $3$ вместо примерно $3{,}68$. & а) 22; б) 3 & а) 22; б) 3,7\\
\addlinespace
\hyperref[task:20]{20} & Высота и медиана в прямоугольном треугольнике & Получен угол $56^\circ$, который является соседним углом при точке на гипотенузе, а не искомым углом между высотой и медианой. & $56^\circ$ & $34^\circ$\\
\bottomrule
\end{longtable}
\normalsize

\newpage
\thispagestyle{reviewfooter}
\vspace*{\fill}
\begin{center}
{\LARGE\bfseries Разбор ошибок}
\end{center}
\vspace*{\fill}

% ---------------- Task 3 ----------------
\newpage
\phantomsection\label{task:3}
\section*{Задание 3}

\subhead{Текст задания}
Прямые $m$ и $n$ параллельны. Найдите угол 3, если угол 1 равен $22^\circ$, а угол 2 равен $72^\circ$.


\begin{center}
\begin{tikzpicture}[scale=0.95,line cap=round,line join=round]
  \draw[line width=0.8pt] (-3,1.2) -- (3,1.2) node[right] {$m$};
  \draw[line width=0.8pt] (-3,-1.2) -- (3,-1.2) node[right] {$n$};
  \draw[line width=0.9pt] (-1.35,2.45) -- (1.35,-2.05);
  \draw[line width=0.9pt] (-2.0,-2.0) -- (1.8,2.0);
  \node at (0.95,1.55) {$22^\circ$};
  \node at (0.18,0.42) {$\angle 3$};
  \node at (0.95,-0.82) {$72^\circ$};
  \draw (-2.5,1.05)--(-2.25,1.35); \draw (-2.2,1.05)--(-1.95,1.35);
  \draw (-2.5,-1.35)--(-2.25,-1.05); \draw (-2.2,-1.35)--(-1.95,-1.05);
\end{tikzpicture}

\small Схема к задаче: оба известных угла нужно перенести к одному развёрнутому углу.
\end{center}

\subhead{Ваше решение}
Записано:
\[
180^\circ-72^\circ=108^\circ.
\]

\subhead{Ваш ответ}
$108^\circ$.

\subhead{Ошибка}
\errorlabel{} пропущен ещё один известный угол, равный $22^\circ$.

\subhead{Где именно возникает ошибка}
Последний корректный факт: угол 2 равен $72^\circ$.

Первый неверный переход --- считать, что искомый угол дополняет только угол $72^\circ$ до $180^\circ$. Под прямой $m$ развёрнутый угол разбит не на две, а на три части: на угол $22^\circ$, на угол 3 и на угол $72^\circ$.

\subhead{Почему это неверно}
Угол, вертикальный углу 1, также равен $22^\circ$. Из параллельности прямых $m$ и $n$ следует, что угол, образованный второй секущей с прямой $m$ в соответствующем положении, равен углу 2, то есть $72^\circ$.

Поэтому под прямой $m$ последовательно расположены три смежные части одного развёрнутого угла. Их сумма обязана быть равна $180^\circ$. Если вычесть только $72^\circ$, часть величиной $22^\circ$ останется неучтённой, поэтому получается завышенный результат.

\subhead{Ключевая идея}
\keyidea{} сначала перенесите оба известных угла к одной точке с помощью вертикальных и соответствующих углов, а затем используйте сумму частей развёрнутого угла.

\subhead{Исправление от последнего правильного шага}
Вместо одного вычитания нужно учесть оба известных угла:
\[
\angle 3=180^\circ-22^\circ-72^\circ=86^\circ.
\]

\subhead{Полное правильное решение}
Так как прямые $m$ и $n$ параллельны, соответствующий угол при верхней прямой, образованный второй секущей, равен углу 2:
\[
72^\circ.
\]

Угол, вертикальный углу 1, равен:
\[
22^\circ.
\]

Под прямой $m$ эти два угла и угол 3 составляют развёрнутый угол, поэтому
\[
22^\circ+\angle 3+72^\circ=180^\circ.
\]
Отсюда
\[
\angle 3=180^\circ-94^\circ=86^\circ.
\]

\subhead{Проверка}
\[
22^\circ+86^\circ+72^\circ=180^\circ,
\]
то есть сумма действительно даёт развёрнутый угол.

\subhead{Итог}
\resultlabel{} правильный ответ: $86^\circ$.

\subhead{Задача на закрепление}
Прямые $a$ и $b$ параллельны. Две секущие пересекаются в точке $P$ на прямой $a$. Угол между первой секущей и прямой $a$ равен $31^\circ$, а соответствующий острый угол второй секущей с прямой $b$ равен $64^\circ$. Найдите угол между секущими, расположенный внутри полосы между параллельными прямыми.

% ---------------- Task 6 ----------------
\newpage
\phantomsection\label{task:6}
\section*{Задание 6}

\subhead{Текст задания}
Для пяти сельскохозяйственных культур известны характеристики урожайности: среднее арифметическое $165$ ц/га, медиана $180$ ц/га, максимум $370$ ц/га, минимум $16$ ц/га. Нужно выбрать верную диаграмму долей и найти урожайность картофеля.

\subhead{Ваше решение}
Верно найдено общее значение урожайности пяти культур:
\[
165\cdot 5=825.
\]
Далее записано выражение вида
\[
16+82+180+x+370=825,
\]
получено $x=177$. В ответе указано: а) диаграмма 3; б) $180$ ц/га.

\subhead{Ваш ответ}
а) 3; б) $180$ ц/га.

\subhead{Ошибка}
\errorlabel{} итоговые ответы верны, но промежуточное вычисление содержит внутреннее противоречие с условием о медиане.

\subhead{Где именно возникает ошибка}
Последний надёжно правильный шаг:
\[
165\cdot5=825.
\]

Первый проблемный шаг --- запись точных значений $82$ и $180$ непосредственно по диаграмме и последующее получение $177$ без проверки медианы. Если принять полученный набор
\[
16,\ 82,\ 177,\ 180,\ 370,
\]
то его медиана равна $177$, а по условию должна быть равна $180$.

\subhead{Почему это неверно}
Диаграмма показывает относительные высоты столбцов, а не подписывает все точные урожайности. Поэтому нельзя сначала приписать столбцу точное число только по его визуальной высоте, а затем считать остальные значения как уже надёжно установленные.

Кроме того, для пяти чисел медиана --- это третье число после упорядочивания. Значит, в правильной диаграмме культура с третьим по высоте столбцом должна иметь урожайность ровно $180$ ц/га.

\subhead{Ключевая идея}
\keyidea{} сначала используйте минимум, максимум и медиану для определения порядка столбцов; среднее арифметическое служит дополнительной проверкой масштаба и долей.

\subhead{Исправление от последнего правильного шага}
Общая урожайность равна
\[
825\text{ ц/га}.
\]
Доли минимального и максимального значений составляют соответственно
\[
\frac{16}{825}\approx1{,}9\%,\qquad \frac{370}{825}\approx44{,}8\%.
\]
Это согласуется с диаграммой 3: один столбец очень мал, а один заметно самый большой. На диаграмме 3 третий по высоте столбец соответствует картофелю. Поскольку медиана равна $180$ ц/га, урожайность картофеля равна $180$ ц/га.

\subhead{Полное правильное решение}
Пять значений имеют среднее арифметическое $165$ ц/га, поэтому их сумма равна
\[
165\cdot5=825\text{ ц/га}.
\]

Минимальное значение $16$ ц/га должно соответствовать самому низкому столбцу, а максимальное $370$ ц/га --- самому высокому. Отношение максимума к минимуму больше $23$, поэтому разница между самым высоким и самым низким столбцами должна быть очень большой. Этому соответствует диаграмма 3.

У пяти чисел медиана --- третье число после расположения по возрастанию. На диаграмме 3 третьим по высоте является столбец картофеля. Следовательно,
\[
\text{урожайность картофеля}=180\text{ ц/га}.
\]

\subhead{Проверка}
В выбранной диаграмме самый низкий столбец соответствует минимуму, самый высокий --- максимуму, а картофель занимает третье место по величине, что согласуется с медианой $180$ ц/га.

\subhead{Итог}
\resultlabel{} а) диаграмма 3; б) $180$ ц/га. Итоговые ответы были записаны верно, но промежуточный расчёт $177$ нельзя оставлять без проверки медианы.

\subhead{Задача на закрепление}
У пяти культур средняя урожайность равна $170$ ц/га, медиана равна $185$ ц/га, минимум равен $20$ ц/га, максимум равен $360$ ц/га. Урожайность одной из оставшихся культур равна $70$ ц/га. Найдите урожайность пятой культуры и расположите все пять значений по возрастанию.

% ---------------- Task 8 ----------------
\newpage
\phantomsection\label{task:8}
\section*{Задание 8}

\subhead{Текст задания}
Для площадей пяти горных систем известны характеристики: среднее арифметическое $922$ тыс. км$^2$, медиана $623$ тыс. км$^2$, максимум $2700$ тыс. км$^2$, минимум $77$ тыс. км$^2$. Требуется выбрать правильную диаграмму долей и найти примерную площадь Иранского нагорья.

\subhead{Ваше решение}
Верно найдено:
\[
922\cdot5=4610.
\]
Далее записано
\[
77+273+623+x+2700=4610,
\]
откуда получено $x=937$. В ответе указано: 1) диаграмма 4; 2) $2700$ тыс. км$^2$.

\subhead{Ваш ответ}
1) 4; 2) $2700$ тыс. км$^2$.

\subhead{Ошибка}
\errorlabel{} выбран номер диаграммы, который противоречит собственным вычислениям.

\subhead{Где именно возникает ошибка}
Последний корректный шаг --- получение набора величин
\[
77,\ 273,\ 623,\ 937,\ 2700.
\]
Он имеет нужные минимум, медиану, максимум и сумму $4610$.

Первый неверный шаг --- запись номера 4. Для такого набора самая большая величина должна занимать примерно $59\%$ всей суммы, а минимальная --- менее $2\%$. Эти соотношения показывает диаграмма 3, а не диаграмма 4.

\subhead{Почему это неверно}
Вычислим доли двух крайних значений:
\[
\frac{2700}{4610}\approx58{,}6\%,\qquad
\frac{77}{4610}\approx1{,}7\%.
\]
Следовательно, на верной круговой диаграмме должен быть сектор, занимающий заметно больше половины круга, и очень маленький сектор около двух процентов. Именно это видно на диаграмме 3.

Кроме того, медианное значение $623$ должно быть третьим по величине. В диаграмме 3 порядок секторов согласуется и с этим условием.

\subhead{Ключевая идея}
\keyidea{} после вычисления чисел обязательно сопоставьте их порядок и примерные доли с рисунком; номер диаграммы должен подтверждаться расчётом, а не записываться отдельно от него.

\subhead{Исправление от последнего правильного шага}
По Вашим же вычислениям:
\[
77<273<623<937<2700.
\]
Максимальное значение $2700$ составляет около $59\%$ суммы, а минимальное $77$ --- около $2\%$. Поэтому выбираем диаграмму 3.

На диаграмме 3 Иранское нагорье соответствует самому большому сектору. Значит, его площадь равна максимальному значению:
\[
2700\text{ тыс. км}^2.
\]

\subhead{Полное правильное решение}
Суммарная площадь пяти систем:
\[
922\cdot5=4610\text{ тыс. км}^2.
\]

Максимум:
\[
2700\text{ тыс. км}^2,
\]
то есть примерно
\[
\frac{2700}{4610}\cdot100\%\approx58{,}6\%.
\]
Минимум:
\[
77\text{ тыс. км}^2,
\]
то есть примерно
\[
\frac{77}{4610}\cdot100\%\approx1{,}7\%.
\]

Среди предложенных рисунков только диаграмма 3 одновременно имеет сектор примерно $59\%$ и очень маленький сектор примерно $2\%$, а третий по величине сектор согласуется с медианой $623$.

Иранское нагорье на этой диаграмме --- самый большой сектор, поэтому его площадь равна
\[
2700\text{ тыс. км}^2.
\]

\subhead{Проверка}
Значение $2700$ является максимумом и действительно соответствует самому большому сектору диаграммы 3.

\subhead{Итог}
\resultlabel{} правильные ответы: 1) 3; 2) $2700$ тыс. км$^2$.

\subhead{Задача на закрепление}
Площади пяти природных областей имеют среднее арифметическое $900$ тыс. км$^2$, медиану $600$ тыс. км$^2$, максимум $2500$ тыс. км$^2$ и минимум $100$ тыс. км$^2$. Площадь одной из двух оставшихся областей равна $300$ тыс. км$^2$. Найдите площадь пятой области и вычислите долю каждой области в общей площади в процентах с точностью до десятых.

% ---------------- Task 10 ----------------
\newpage
\phantomsection\label{task:10}
\section*{Задание 10}

\subhead{Текст задания}
Укажите номера верных утверждений:
\begin{enumerate}
\item Если расстояние между центрами двух окружностей равно сумме их диаметров, то окружности касаются.
\item Центр окружности, описанной около треугольника со сторонами $3$, $4$, $5$, находится вне этого треугольника.
\item Если радиус окружности равен $3$, а расстояние от центра окружности до прямой равно $2$, то прямая и окружность пересекаются.
\item Для точки, лежащей на окружности, расстояние до центра окружности равно радиусу.
\end{enumerate}


\begin{center}
\begin{tikzpicture}[scale=0.82,line cap=round,line join=round]
  \begin{scope}[xshift=-4.1cm]
    \draw (0,0) circle (0.72);
    \draw (3.0,0) circle (0.78);
    \fill (0,0) circle (1.4pt) node[below=3pt] {$O_1$};
    \fill (3.0,0) circle (1.4pt) node[below=3pt] {$O_2$};
    \draw[<->] (0,1.18) -- (3.0,1.18) node[midway,fill=white,inner sep=1pt] {$d$};
    \node[align=center,text width=3.4cm] at (1.5,-1.52) {$d=2r_1+2r_2$\\[-1pt]\small окружности не касаются};
  \end{scope}
  \begin{scope}[xshift=3.9cm]
    \draw (0,0) circle (1.25);
    \fill (0,0) circle (1.4pt) node[left=3pt] {$O$};
    \draw (-2,-0.78) -- (2,-0.78);
    \draw[dashed] (0,0) -- (0,-0.78) node[midway,right] {$2$};
    \draw (0,0) -- (0.88,0.88) node[midway,above left] {$3$};
    \node[align=center,text width=3.8cm] at (0,-1.58) {$2<3$\\[-1pt]\small прямая пересекает окружность};
  \end{scope}
\end{tikzpicture}

\small Наглядная проверка утверждений 1 и 3.
\end{center}

\subhead{Ваше решение}
Отмечено: 1 --- верно; 2 --- неверно; 3 --- верно; 4 --- верно.

\subhead{Ваш ответ}
Верными признаны утверждения 1, 3 и 4.

\subhead{Ошибка}
\errorlabel{} утверждение 1 ошибочно признано верным.

\subhead{Где именно возникает ошибка}
Первый неверный шаг --- использование суммы диаметров как условия касания двух окружностей.

\subhead{Почему это неверно}
При внешнем касании расстояние между центрами равно сумме \emph{радиусов}:
\[
d=r_1+r_2.
\]
Диаметры равны $2r_1$ и $2r_2$, поэтому сумма диаметров равна
\[
2r_1+2r_2=2(r_1+r_2),
\]
то есть в два раза больше расстояния, необходимого для внешнего касания. Поэтому утверждение 1 неверно.

Утверждение 2 также неверно: треугольник со сторонами $3$, $4$, $5$ прямоугольный, а центр описанной окружности прямоугольного треугольника находится в середине гипотенузы, то есть на стороне треугольника, а не вне него.

Утверждение 3 верно, потому что расстояние от центра до прямой меньше радиуса: $2<3$, значит прямая пересекает окружность в двух точках.

Утверждение 4 верно по определению окружности: все её точки находятся от центра на расстоянии, равном радиусу.

\subhead{Ключевая идея}
\keyidea{} для окружностей различайте радиус и диаметр, а положение прямой относительно окружности определяйте сравнением расстояния от центра до прямой с радиусом.

\subhead{Исправление от последнего правильного шага}
Сохраняем оценки утверждений 2, 3 и 4, но меняем оценку утверждения 1 на «неверно».

\subhead{Полное правильное решение}
\begin{enumerate}
\item Неверно: для касания нужно $d=r_1+r_2$, а не сумма диаметров.
\item Неверно: треугольник $3$--$4$--$5$ прямоугольный; центр описанной окружности лежит в середине гипотенузы.
\item Верно: $2<3$, поэтому прямая пересекает окружность.
\item Верно: это определяющее свойство точки окружности.
\end{enumerate}
Следовательно, верны утверждения 3 и 4.

\subhead{Проверка}
Для утверждения 3 при расстоянии, равном радиусу, было бы касание, при расстоянии больше радиуса пересечения не было бы. Здесь расстояние меньше радиуса, поэтому пересечение действительно есть.

\subhead{Итог}
\resultlabel{} правильный ответ: 3, 4.

\subhead{Задача на закрепление}
Укажите номера верных утверждений: 1) если расстояние между центрами двух окружностей равно сумме их радиусов, окружности имеют внешнее касание; 2) центр описанной окружности равностороннего треугольника лежит вне треугольника; 3) если радиус окружности равен $5$, а расстояние от центра до прямой равно $3$, прямая пересекает окружность; 4) для любой точки внутри окружности расстояние до центра равно радиусу.

% ---------------- Task 11 ----------------
\newpage
\phantomsection\label{task:11}
\section*{Задание 11}

\subhead{Текст задания}
Даны схема дорог и таблица длин дорог между пунктами $\text{П1},\ldots,\text{П7}$. Нумерация в таблице не совпадает с буквенными обозначениями на схеме. Требуется определить длину дороги из пункта Г в пункт Е.


\begin{center}
\begin{tikzpicture}[scale=0.95,line cap=round,line join=round]
  \coordinate (A) at (-3.0,1.2);
  \coordinate (V) at (-1.8,0.2);
  \coordinate (D) at (-3.35,-1.0);
  \coordinate (G) at (0.15,0.35);
  \coordinate (B) at (1.15,1.55);
  \coordinate (E) at (1.15,-1.15);
  \coordinate (K) at (3.25,-1.15);
  \draw[line width=0.9pt] (A)--(V)--(D);
  \draw[line width=0.9pt] (V)--(G)--(B);
  \draw[line width=0.9pt] (G)--(E)--(K);
  \foreach \P/\lab/\pos in {A/А/above left,V/В/above,D/Д/below left,G/Г/above,E/Е/below,B/Б/above,K/К/below} {
    \fill (\P) circle (2pt) node[\pos=3pt] {\lab};
  }
  \node[align=center] at (0,-2.05) {По структуре: П3 $\leftrightarrow$ В,\quad
  П4 $\leftrightarrow$ Г,\quad
  П6 $\leftrightarrow$ Е.};
\end{tikzpicture}

\small Схема графа: длины нарисованных отрезков значения не имеют; важны связи между вершинами.
\end{center}

\subhead{Ваше решение}
Записано: «не знаю».

\subhead{Ваш ответ}
Не указан числовой ответ.

\subhead{Ошибка}
\errorlabel{} решение не выполнено: не сопоставлены вершины графа с пунктами таблицы.

\subhead{Где именно возникает ошибка}
В задаче можно начать с количества дорог, выходящих из каждой вершины. Этот шаг отсутствует, поэтому сопоставление не построено.

\subhead{Почему это неверно}
Форма рисунка и длины нарисованных отрезков не используются. Надёжный признак --- степень вершины, то есть количество дорог, которые к ней подходят.

В таблице:
\begin{itemize}
\item П3 соединён с П1, П2 и П4 --- степень 3;
\item П4 соединён с П3, П5 и П6 --- степень 3;
\item П6 соединён с П4 и П7 --- степень 2.
\end{itemize}

На схеме:
\begin{itemize}
\item пункт В соединён с двумя тупиковыми пунктами и с Г --- степень 3;
\item пункт Г соединён с В, одним тупиковым пунктом и Е --- степень 3;
\item пункт Е соединён с Г и тупиковым пунктом К --- степень 2.
\end{itemize}

По структуре цепочки получаем соответствие П3 $\leftrightarrow$ В, П4 $\leftrightarrow$ Г, П6 $\leftrightarrow$ Е.

\subhead{Ключевая идея}
\keyidea{} сначала сопоставляйте вершины по числу соседей и по типу соседей, а длину дороги считывайте из таблицы только после установления соответствия.

\subhead{Исправление от последнего правильного шага}
Нужно определить, какие номера соответствуют Г и Е. Пункт Е имеет степень 2 и соединён с вершиной степени 3 и тупиком; в таблице такую роль играет П6. Пункт Г --- сосед Е степени 3, значит ему соответствует П4.

\subhead{Полное правильное решение}
По таблице рёбра имеют вид:
\[
\text{П1--П3},\quad \text{П2--П3},\quad \text{П3--П4},\quad
\text{П4--П5},\quad \text{П4--П6},\quad \text{П6--П7}.
\]

Вершины П3 и П4 имеют степень 3, П6 --- степень 2. На буквенной схеме вершины В и Г имеют степень 3, а Е --- степень 2.

При этом вершина Е степени 2 соединена с Г степени 3 и с тупиком К. В таблице П6 соединена с П4 степени 3 и с тупиком П7. Следовательно,
\[
\text{Е}\leftrightarrow\text{П6},\qquad
\text{Г}\leftrightarrow\text{П4}.
\]

В таблице длина дороги П4--П6 равна $12$ км. Поэтому длина дороги Г--Е равна $12$ км.

\subhead{Проверка}
После этого сопоставления оставшаяся структура графа также совпадает: П3 соответствует В, а четыре вершины степени 1 соответствуют четырём тупиковым пунктам схемы.

\subhead{Итог}
\resultlabel{} длина дороги из Г в Е равна $12$ км.

\subhead{Задача на закрепление}
В схеме дорог вершина $X$ соединена с $A$, $B$ и $Y$; вершина $Y$ соединена с $X$, $C$ и $Z$; вершина $Z$ соединена с $Y$ и $D$. В таблице пункт П2 соединён с П1 дорогой длиной $7$ км, с П3 дорогой длиной $11$ км и с П5 дорогой длиной $9$ км; П5 соединён также с П4 дорогой длиной $13$ км и с П6 дорогой длиной $15$ км; П6 соединён также с П7 дорогой длиной $20$ км. Найдите длину дороги $Y$--$Z$.

% ---------------- Task 16 ----------------
\newpage
\phantomsection\label{task:16}
\section*{Задание 16}

\subhead{Текст задания}
Купили $3$ кг яблок по $170$ руб./кг, $1$ кг изюма по $210$ руб./кг и $1$ кг кураги по $590$ руб./кг. 1) Найдите среднюю стоимость $1$ кг смеси. 2) После добавления чернослива по $900$ руб./кг средняя стоимость стала $350$ руб./кг. Сколько килограммов чернослива купили?

\subhead{Ваше решение}
В первой части стоимость смеси посчитана как стоимость трёх килограммов яблок, одного килограмма изюма и одного килограмма кураги, после чего получено $262$ руб./кг.

Во второй части начата запись новой средней стоимости, но уравнение с неизвестной массой чернослива не доведено до решения.

\subhead{Ваш ответ}
1) $262$ руб.; 2) ответ не завершён.

\subhead{Ошибка}
\errorlabel{} вторая часть решения остановлена до составления правильного уравнения для массы и стоимости смеси.

\subhead{Где именно возникает ошибка}
Первая часть выполнена правильно. Последний корректный результат:
\[
\text{стоимость исходной смеси}=1310\text{ руб.},\qquad
\text{масса исходной смеси}=5\text{ кг}.
\]

Далее нужно было обозначить массу чернослива переменной и одновременно изменить и общую стоимость, и общую массу.

\subhead{Почему это неверно}
Средняя стоимость одного килограмма --- это отношение общей стоимости смеси к общей массе. После добавления $x$ килограммов чернослива:
\begin{itemize}
\item стоимость увеличивается на $900x$ рублей;
\item масса увеличивается на $x$ килограммов.
\end{itemize}
Нельзя добавлять цену одного килограмма как одно фиксированное число, если неизвестно, сколько килограммов куплено.

\subhead{Ключевая идея}
\keyidea{} при изменении смеси записывайте отдельно новую общую стоимость и новую общую массу, затем приравнивайте их отношение к новой средней цене.

\subhead{Исправление от последнего правильного шага}
Пусть купили $x$ кг чернослива. Тогда
\[
\frac{1310+900x}{5+x}=350.
\]
Решаем это уравнение.

\subhead{Полное правильное решение}
Стоимость яблок:
\[
3\cdot170=510\text{ руб.}
\]
Общая стоимость исходной смеси:
\[
510+210+590=1310\text{ руб.}
\]
Общая масса:
\[
3+1+1=5\text{ кг}.
\]
Поэтому средняя стоимость одного килограмма:
\[
\frac{1310}{5}=262\text{ руб./кг}.
\]

Для второй части обозначим массу чернослива через $x$ кг. После покупки общая стоимость будет
\[
1310+900x,
\]
а общая масса
\[
5+x.
\]
По условию новая средняя стоимость равна $350$ руб./кг:
\[
\frac{1310+900x}{5+x}=350.
\]
Умножаем обе части на $5+x$:
\[
1310+900x=1750+350x.
\]
Переносим слагаемые:
\[
550x=440.
\]
Отсюда
\[
x=0{,}8.
\]

\subhead{Проверка}
При $x=0{,}8$ общая стоимость равна
\[
1310+900\cdot0{,}8=2030\text{ руб.},
\]
а масса
\[
5+0{,}8=5{,}8\text{ кг}.
\]
Проверяем:
\[
\frac{2030}{5{,}8}=350\text{ руб./кг}.
\]

\subhead{Итог}
\resultlabel{} 1) $262$ руб./кг; 2) $0{,}8$ кг чернослива.

\subhead{Задача на закрепление}
Для смеси взяли $2$ кг яблок по $150$ руб./кг, $1$ кг изюма по $200$ руб./кг и $1$ кг кураги по $500$ руб./кг. Затем добавили чернослив по $900$ руб./кг, после чего средняя стоимость смеси стала $400$ руб./кг. Сколько килограммов чернослива добавили?

% ---------------- Task 18 ----------------
\newpage
\phantomsection\label{task:18}
\section*{Задание 18}

\subhead{Текст задания}
По диаграмме результатов контрольной работы: оценку «2» получили 3 человека, «3» --- 6 человек, «4» --- 8 человек, «5» --- 5 человек. 1) Сколько человек писали работу? 2) Найдите среднюю оценку и округлите до десятых.

\subhead{Ваше решение}
Верно найдено число учеников:
\[
3+6+8+5=22.
\]
Также верно получена сумма всех оценок:
\[
81.
\]
Далее записано:
\[
\frac{81}{22}=3.
\]

\subhead{Ваш ответ}
а) 22; б) 3.

\subhead{Ошибка}
\errorlabel{} неверно выполнено деление $81$ на $22$ и не проведено требуемое округление до десятых.

\subhead{Где именно возникает ошибка}
Последний правильный шаг --- сумма оценок $81$ и количество учеников $22$.

Первый неверный шаг:
\[
\frac{81}{22}=3.
\]
На самом деле $22\cdot3=66$, поэтому частное должно быть заметно больше трёх.

\subhead{Почему это неверно}
Средняя оценка определяется как сумма всех оценок, делённая на число учеников. Так как
\[
81-66=15,
\]
после целой части $3$ остаётся ещё дробная часть
\[
\frac{15}{22}\approx0{,}68.
\]
Следовательно, среднее значение примерно равно $3{,}68$, а после округления до десятых --- $3{,}7$.

\subhead{Ключевая идея}
\keyidea{} после получения взвешенной суммы не забывайте делить с дробной частью и выполнять именно то округление, которое требует условие.

\subhead{Исправление от последнего правильного шага}
Из уже найденных чисел:
\[
\overline{x}=\frac{81}{22}\approx3{,}6818.
\]
До десятых:
\[
3{,}7.
\]

\subhead{Полное правильное решение}
Всего работу писали:
\[
3+6+8+5=22\text{ человека}.
\]

Сумма всех выставленных оценок:
\[
2\cdot3+3\cdot6+4\cdot8+5\cdot5.
\]
Вычисляем:
\[
6+18+32+25=81.
\]
Средняя оценка:
\[
\frac{81}{22}\approx3{,}6818.
\]
Округляем до десятых. В сотых стоит цифра $8$, поэтому цифру десятых увеличиваем на единицу:
\[
3{,}7.
\]

\subhead{Проверка}
Среднее значение должно лежать между минимальной и максимальной оценками, то есть между $2$ и $5$. Значение $3{,}7$ этому условию соответствует и находится ближе к $4$, что согласуется с диаграммой: оценка «4» встречается чаще всего.

\subhead{Итог}
\resultlabel{} а) 22 человека; б) средняя оценка $3{,}7$.

\subhead{Задача на закрепление}
На другой контрольной работе оценку «2» получили 2 человека, «3» --- 5 человек, «4» --- 7 человек, «5» --- 6 человек. Найдите среднюю оценку и округлите результат до десятых.

% ---------------- Task 20 ----------------
\newpage
\phantomsection\label{task:20}
\section*{Задание 20}

\subhead{Текст задания}
Острые углы прямоугольного треугольника равны $62^\circ$ и $28^\circ$. Найдите угол между высотой и медианой, проведёнными из вершины прямого угла.


\begin{center}
\begin{tikzpicture}[scale=0.95,line cap=round,line join=round]
  \coordinate (A) at (0,0);
  \coordinate (B) at (6,0);
  \coordinate (C) at (4.676,2.487);
  \coordinate (F) at (3,0);
  \coordinate (H) at (4.676,0);
  \draw[line width=0.9pt] (A)--(C)--(B)--cycle;
  \draw[line width=0.9pt] (C)--(F);
  \draw[line width=0.9pt] (C)--(H);
  \node[below left] at (A) {$A$};
  \node[below right] at (B) {$B$};
  \node[above] at (C) {$C$};
  \node[below] at (F) {$F$};
  \node[below] at (H) {$H$};
  \draw (H) ++(-0.25,0) -- ++(0,0.25) -- ++(0.25,0);
  \draw (1.47,-0.10)--(1.53,0.10);
  \draw (4.47,-0.10)--(4.53,0.10);
  \node at (0.73,0.34) {$28^\circ$};
  \node at (5.35,0.36) {$62^\circ$};
  \node at (4.29,1.55) {$34^\circ$};
  \node[align=center] at (3,-0.82) {$CF$ --- медиана,\qquad $CH$ --- высота};
\end{tikzpicture}

\small Искомый угол расположен при вершине $C$, между лучами $CF$ и $CH$.
\end{center}

\subhead{Ваше решение}
В рисунке обозначены медиана $CF$ и высота $CH$. Верно использовано свойство середины гипотенузы и получен угол при точке $F$, равный $56^\circ$. В конце записан ответ $56^\circ$.

\subhead{Ваш ответ}
$56^\circ$.

\subhead{Ошибка}
\errorlabel{} найден угол между медианой и гипотенузой, а требуется угол между медианой и высотой.

\subhead{Где именно возникает ошибка}
Последний правильный шаг --- получение
\[
\angle CFH=56^\circ.
\]

Первый неверный шаг --- принять этот угол за искомый. Искомый угол находится при вершине $C$, между лучами $CF$ и $CH$.

\subhead{Почему это неверно}
Так как $CH$ --- высота к гипотенузе $AB$, то
\[
CH\perp AB.
\]
Следовательно, в треугольнике $CFH$ угол при $H$ равен $90^\circ$. Если угол при $F$ равен $56^\circ$, то третий угол при $C$ нельзя также принять равным $56^\circ$: сумма углов треугольника должна быть $180^\circ$.

\subhead{Ключевая идея}
\keyidea{} после нахождения вспомогательного угла обязательно проверьте, при какой вершине он расположен. Искомый угол должен иметь стороны, совпадающие с высотой и медианой.

\subhead{Исправление от последнего правильного шага}
Из треугольника $CFH$:
\[
\angle FCH=180^\circ-90^\circ-56^\circ=34^\circ.
\]

\subhead{Полное правильное решение}
Пусть $C$ --- вершина прямого угла, $F$ --- середина гипотенузы $AB$, поэтому $CF$ --- медиана, а $H$ --- основание высоты $CH$.

В прямоугольном треугольнике середина гипотенузы равноудалена от всех трёх вершин, поэтому
\[
AF=CF.
\]
Треугольник $ACF$ равнобедренный. Угол при $A$ равен $28^\circ$, следовательно,
\[
\angle ACF=28^\circ.
\]

В прямоугольном треугольнике $ACH$ угол при $H$ равен $90^\circ$, а угол при $A$ равен $28^\circ$. Поэтому
\[
\angle ACH=180^\circ-90^\circ-28^\circ=62^\circ.
\]

Луч $CF$ расположен внутри угла $ACH$, значит угол между медианой и высотой равен
\[
\angle FCH=62^\circ-28^\circ=34^\circ.
\]

\subhead{Проверка}
Другой способ: в треугольнике $ACF$
\[
\angle AFC=180^\circ-28^\circ-28^\circ=124^\circ.
\]
Так как $A$, $F$, $H$ лежат на одной прямой,
\[
\angle CFH=180^\circ-124^\circ=56^\circ.
\]
Тогда
\[
\angle FCH=180^\circ-90^\circ-56^\circ=34^\circ.
\]
Получен тот же результат.

\subhead{Итог}
\resultlabel{} угол между высотой и медианой равен $34^\circ$.

\subhead{Задача на закрепление}
Острые углы прямоугольного треугольника равны $65^\circ$ и $25^\circ$. Найдите угол между высотой и медианой, проведёнными из вершины прямого угла.

\end{document}